\scanpage{10}
\begin{center}- 43 -\end{center}
\hfill n$^\circ$ 387

le fait que $A$ est ouvert pour $T_K$, montrer que $V$ est un $A$-module de type fini et conclure à l'aide du § 3, n$^\circ$ 5, prop. 9. Pour la réciproque, on observera que $A$ est un anneau local dont l'idéal maximal $\underline m$ et $(0)$ sont les seuls idéaux premiers, et dans lequel tout idéal $\underline a\ne0$ est par suite $\underline m$-primaire). La clôture intégrale de $A$ est alors $V$. Donner un exemple où $A\ne V$ (prendre pour $V$ un anneau de $A$ est alors $V$. Donner un exemple où $A\ne V$ (prendre pour $V$ un anneau de séries formelles $k[[T]]$, où $k$ est un corps).

c) En déduire un exemple de corps topologique complet non discret, dont la topologie n'est pas localement rétrobornée (utiliser b) et l'exerc. 2).

4) Soient $v$ une valuation sur un corps $K$, $A$ son anneau, $\underline m$ son idéal maximal.

a) Pour que la topologie définie par $v$ sur $A$ soit identique à la topologie $\underline m$-adique, il faut et il suffit que $A$ soit un corps ou un anneau de valuation discrète.

b) Pour que la topologie définie par $v$ fasse de $A$ un anneau strictement linéairement compact (chap. III, § 2, exerc. 21), il faut et il suffit que $A$ soit un corps ou un anneau de valuation discrète (utiliser a) et le chap. III, § 2, exerc. 22 a)).

\textsuperscript{*}5) a) Soient $K$ un corps, $v$ une valuation sur $K$, $\Gamma$ le groupe des ordres de $v$. Pour que $K$ soit linéairement compact (chap. III, § 2, exerc. 15) pour la topologie $T_v$, il faut et il suffit que pour toute partie \uline{bien ordonnée} $B$ de $\Gamma$ et toute famille $(a_\beta)_{\beta\in B}$ d'éléments de $K$ telle que, pour $\lambda<\mu<\nu$, on ait $v(a_\lambda-a_\mu)<v(a_\mu-a_\nu)$, il existe un élément $a\in K$ tel que $v(a-a_\lambda)=v(a_\lambda-a_\mu)$ pour tout couple d'indices $\lambda,\mu$ tels que $\lambda<\mu$. (Utiliser l'exerc. de \uline{Ens.}, chap. III, 2$^{\mathrm e}$ éd., § 2). Lorsqu'il en est ainsi, l'anneau $A$ de la valuation $v$ est aussi linéairement compact pour la topologie discrète.

b) Montrer que le corps valué $S(\Gamma,k)=K$ défini au § 3, exerc. 2, est linéairement compact pour $T_v$. On prend pour $\Gamma$ le groupe totalement ordonné $\underline Q$ des nombres rationnels, et on considère le sous-anneau $K_{\unsure{o|0}}$ de $K$ formé des $x=(x_\alpha)$ tels que l'ensemble des $\alpha\in\underline Q$ pour lesquels $x_\alpha\ne0$ soit fini ou soit l'ensemble des points d'une suite strictement croissante tendant vers $+\infty$. Montrer que $K_{\unsure{o|0}}$ est un corps qui est complet pour la valuation $v_{\unsure{o|0}}$ induite par la valuation canonique $v$ de $K$, et que $v_{\unsure{o|0}}$ a même groupe des valeurs et même corps résiduel que $v$,
\mspage{9}
